\documentclass[10pt]{article}
\usepackage[margin=20mm]{geometry}
\usepackage{amsmath,amssymb,mathtools,fontspec,unicode-math,xurl,hyperref}
\IfFontExistsTF{DejaVu Serif}{\setmainfont{DejaVu Serif}}{\setmainfont{Latin Modern Roman}}
\setmathfont{latinmodern-math.otf}
\setlength{\parindent}{0pt}\setlength{\parskip}{6pt}\sloppy
\title{Technical review of
Multilinear Hyperquiver Representations}
\author{AI-assisted technical review}
\date{}
\begin{document}
\maketitle
Abstract: The general dimension and emptiness assertions admit an explicit counterexample using repeated vertices and permuted copies of one generic tensor. The stated Bertini theorem also admits counterexamples to emptiness and smoothness. Three additional local assertions concerning global sections, homological fixed classes, and convergence fail in elementary examples. The inverse problem’s Hermitian-adjoint convention and exact-arithmetic recurrence were checked without establishing a defect. The overall supported impact is Major, confined to the claims and source version specified below; scientific adjudication and independent acceptance remain pending.\par
\section*{Source and scope}
The reviewed paper is by Tommi Muller, Vidit Nanda and Anna Seigal, DOI 10.1007/s10208-025-09692-z (https://doi.org/10.1007/s10208-025-09692-z). The source is the author-hosted, publisher-typeset online-first PDF associated with 14 February 2025, comprising 43 PDF pages. All page references below mean one-based PDF pages. Equality with the later repaginated issue file has not been established.\par
The focused review covers Definitions 2.1–2.6, Theorems 1.1 and 3.1, the section map in Proposition 4.5, Definition 5.2 and Theorem 5.3, the relevant genericity argument in Section 6, Examples 3.12–3.13, and the inverse problem in Section 8. Standard special-case coefficient calculations were also compared with retained exact computations. This is not a verification of every example, external reference, Appendix A result, or floating-point implementation. No relevant author implementation or separate supplement was located in the preserved acquisition search; absence of code is not a mathematical defect.\par
Throughout, tensor contraction uses the bilinear pairing \(u^{\mathsf T}v\), without conjugation (Remark 2.4, p. 8). Thus nonzero isotropic vectors are permitted. Definitions 2.1–2.5 explicitly allow repeated vertices and related tensors on distinct edges. These features are essential to the first problem.\par
\section*{A counterexample to the main theorem}
Problem 1 (Generic representations can violate the asserted expected dimension).  Theorem 3.1 (p. 9) asserts that a generic representation has no singular vector tuple if \[N=\sum_{v\in V}(d_v-1)-\sum_{e\in E}(d_{t(e)}-1)<0,\] and otherwise, when nonempty, has pure dimension \(N\) and the stated Chern-class degree. The following representation satisfies its stated genericity conditions and contradicts both branches.\par
Equation reference: eq:contractions. Take two vertices with spaces \(\mathbb C^2\) and \(\mathbb C^q\), where \(q\geq2\). There are three edges, each with ordered vertex tuple \((1,1,1,2)\) and target vertex 1. Choose one generic tensor \[T\in\mathbb C^2\otimes\mathbb C^2\otimes\mathbb C^2\otimes\mathbb C^q.\] Assign its three permuted copies so that the edge contractions are \[
c_0=T(\,\cdot,x,x,y),\quad
c_1=T(x,\,\cdot,x,y),\quad
c_2=T(x,x,\,\cdot,y).\] All three edges form one block of the partition in Definition 2.5. The reference slot orders are \((0,1,2,3)\), \((1,0,2,3)\) and \((2,1,0,3)\). Their pairwise induced permutations send the output slots to three distinct positions, satisfying condition (i.c); the preserved exact check verifies all 18 ordered comparisons. For a generic \(T\), the only tensor constraint is precisely the prescribed permutation relation. No symmetry of \(T\) is imposed. Repeated occurrences of vertex 1 are allowed by Remark 2.2.\par
Set \(x=(1,i)^{\mathsf T}\) and define the linear functional \[\ell_T(y)=T(x,x,x,y).\] Because \(x^{\mathsf T}x=0\) in \(\mathbb C^2\), the orthogonal complement \(x^\perp\) equals \(\operatorname{span}(x)\). Each contraction in [eq:contractions] (\#eq:contractions) satisfies \[x^{\mathsf T}c_j=\ell_T(y).\] For every \(T\), the kernel of \(\ell_T:\mathbb C^q\to\mathbb C\) contains a nonzero vector. Any such \(y\) therefore gives \(c_j\in\operatorname{span}(x)\) for all three edges. Consequently \(([x],[y])\) is a singular vector tuple in the exact sense of Definition 2.6.\par
For generic \(T\), \(\ell_T\) is nonzero: failure is a proper linear condition on \(T\), as a tensor with its \((0,0,0,0)\) entry equal to one already shows. Thus these solutions contain a projective fibre \(\mathbb P^{q-2}\), whereas the theorem predicts \[N=(2-1)+(q-1)-3(2-1)=q-3.\] For \(q=2\), \(N=-1\) but the variety is nonempty for every \(T\); this also contradicts the equal-dimension statement in Theorem 1.1 (p. 5). For \(q=3\), \(N=0\) but the variety contains a projective line. The formal coefficient in (3.1) is then \[[h_1h_2^2](3h_1+h_2)^3=9,\] which cannot count the positive-dimensional zero set. The same construction uses \(x=(1,-i)^{\mathsf T}\).\par
The argument is symbolic and uniform in \(T\), so it does not rely on calling one numerical tensor generic. The retained numerical-algebraic witness is only a corroborating specialization: for lexicographically ordered tensor entries \(2,3,5,\ldots,53\), \(x=(1,i)^{\mathsf T}\) and \[y=(96+4i,-87-8i)^{\mathsf T},\] the three outputs are \[(48-136i,136+48i),\quad
(350-288i,288+350i),\quad
(-64-226i,226-64i),\] all proportional to \(x\). The symbolic evaluation map has rank three at the checked nonisotropic point and rank one at the isotropic point. This is the relevant rank collapse, not an OCR or transpose convention.\par
\subsection*{Impact and counterevidence}
The failure concerns the paper’s central generalization to permitted tensor networks and supports Major severity. It does not refute the classical generic single-tensor counts or every restricted hyperquiver case. Imposing independence of all edge tensors would change this construction’s hypotheses, but Definition 2.5 expressly permits the shared tensor block. Excluding isotropic tuples would change the defined solution space. No complete corrected general theorem is asserted here.\par
\section*{The stated Bertini hypotheses are insufficient}
Problem 2 (Almost generation as defined does not ensure emptiness or smoothness).  Definition 5.2 and Theorem 5.3 (pp. 22–24) allow exceptional subvarieties with lower-rank generated subbundles, without the rank-versus-dimension condition needed to exclude their zeros. Two examples satisfy the stated definition.\par
For the emptiness assertion, take \[X=\mathbb P^1,\qquad B=\mathcal O(1)\oplus\mathcal O(1),\qquad
\Lambda=\{(aX_0,bX_0):a,b\in\mathbb C\}.\] Use the exceptional subvariety \(Y=\{[0:1]\}\) and the zero subbundle \(B_Y=0\). Off \(Y\), evaluation of \(\Lambda\) spans the rank-two fibre. Along \(Y\), every section is zero and hence generates exactly the stipulated zero subbundle. The exceptional locus is proper and smooth; there are no unresolved intersection conditions. Nevertheless, every nonzero section vanishes at \(Y\), although \(\operatorname{rank}B=2>\dim X=1\). This contradicts Theorem 5.3(a). The proof’s incidence object has fibre dimension zero off \(Y\) and two over \(Y\), so its asserted vector-bundle structure (p. 23) does not hold.\par
For smoothness, take \[X=\mathbb P^2,\qquad B=\mathcal O(2),\qquad
\Lambda=\operatorname{span}\{X_0X_1\}.\] Use the lines \(Y_1=\{X_0=0\}\), \(Y_2=\{X_1=0\}\) and their intersection \(Y_3\), with zero subbundles on all three strata. Evaluation is onto away from the two lines and zero on the exceptional loci. Every nonzero, hence generic, section vanishes on their union. At \([0:0:1]\), its local equation is \(uv=0\), whose gradient vanishes and whose tangent space has dimension two rather than the curve’s dimension one. Thus the smoothness assertion in Theorem 5.3(c) fails.\par
\subsection*{Counterevidence and bounded correction}
The consulted Friedland–Ottaviani version [FO], Definition 5 and Theorem 6 (PDF pp. 10–11), includes a requirement that the generated subbundle over each exceptional locus have rank strictly greater than that locus’s dimension. Neither example meets that stronger requirement. The present counterexamples therefore concern the weaker statement in the reviewed paper, not that cited result. Restoring a suitable inequality explains how these two examples can be excluded; it does not, by itself, verify all hypotheses for the hyperquiver construction. The first problem shows why a blanket repair of the main theorem cannot be inferred. Because Theorem 5.3 is a central proof ingredient and is itself false as stated, the supported impact is Major.\par
\section*{Additional local errors}
Problem 3 (The map in Proposition 4.5 need not be bijective).  Proposition 4.5 (pp. 21–22) identifies the tensor parameter space with all global sections of \(B(R)\). The allowed Jordan quiver already contradicts injectivity.\par
For one vertex of dimension two and one loop, \(B(R)\simeq\mathcal O(1)\otimes Q\simeq\mathcal O(2)\) on \(\mathbb P^1\). Its global-section space has dimension three, whereas the matrix space has dimension four. For \(A=\left(\begin{smallmatrix}a&b\\c&d\end{smallmatrix}\right)\) the induced section is represented by \[\det[x,Ax]=cx_0^2+(d-a)x_0x_1-bx_1^2.\] Its kernel is the scalar matrices. The proof chooses a target vector independently of the source vectors; for a self-loop these are the same projective variable.\par
There is also a surjectivity obstruction without a loop. For two two-dimensional vertices and edges labeled by \(A\) and \(A^{\mathsf T}\) in a single tensor block, the parameter space has dimension four, while \[B(R)\simeq \mathcal O(1,1)\oplus\mathcal O(1,1)\] has eight independent global sections. The linked assignments cannot parametrize that entire space.\par
The construction of \(L\) and the zero-locus equivalence in Proposition 4.3 remain meaningful. Working with \(\operatorname{im}L\) is the natural local correction, and the subsequent almost-generation discussion explicitly uses an image linear system. This report does not attribute a separate failure of the final count solely to bijectivity. The supported local severity is Minor; the demonstrated core failures are already isolated above.\par
Problem 4 (Connected parameter spaces do not give a unique fixed homology class).  Example 3.12 (p. 17) states that a connected parameter space gives exactly one nonzero homology class fixed up to scaling. Let \(P\) be a point and \(X=S^1\vee S^1\), and let \(f\) map the two circles to themselves with degrees two and three. These are compact triangulable spaces and a continuous map of the prescribed kind. On \(H_1(X;\mathbb C)\), \[T_1=f_*=\begin{pmatrix}2&0\\0&3\end{pmatrix}.\] The definition in the example permits an arbitrary scalar eigenvalue. It therefore yields two distinct projective eigenlines, not one. If \(f\) is instead the identity, all of \(\mathbb P^1\) consists of fixed classes, which also contradicts the example’s unconditional dimension \(d-1=0\).\par
Requiring \(\lambda=1\) would change the stated definition and does not recover the claimed count in the first example. An additional genericity assumption is also insufficient to imply uniqueness: a generic two-dimensional matrix has two distinct eigenlines. This is a substantive local application error, graded Minor; no dependence of the main tensor theorem on it is demonstrated.\par
Problem 5 (The message-passing update need not converge).  Example 3.13, equation (3.8) and the next sentence (p. 18), assert convergence to a fixed point without a contraction or spectral hypothesis. In the allowed one-vertex, one-loop case, set \[A=\operatorname{diag}(2,-2),\qquad x^{(0)}=(1,1)^{\mathsf T}.\] Then \(x^{(k)}=2^k(1,(-1)^k)^{\mathsf T}\). The vector norms diverge, and the projective directions alternate between \([1:1]\) and \([1:-1]\). Normalizing the norm does not remove this period-two behaviour.\par
The characterization of invariant directions remains useful; it is not a convergence theorem. A dominant-eigenvalue or contraction hypothesis would be additional and must be specified before a convergence statement is justified. This local dynamical claim is graded Minor.\par
\section*{Checks that did not establish an issue}
For the inverse problem (pp. 32–38), the source images confirm conjugated input vectors in the adjoint contractions of equation (8.2) and Algorithm 1, lines 8 and 19. These conjugates are required by the Hermitian inner product used for the minimum-norm problem and are present in the actual source. A supposed missing-conjugation error is therefore not supported.\par
The affine solution expression in Theorem 8.1 was checked against the contraction constraint. Algorithm 1 is an exact-arithmetic conjugate-gradient recurrence for the corresponding normal equation; the orthogonality and error-invariant arguments support its stated finite-dimensional alternative in the inspected scope. The retained exact complex-arithmetic examples give a consistent solution after two iterations and detect inconsistency in a second case. Those two finite tests corroborate, but do not constitute, the general argument. No floating-point stability claim or execution of an author implementation is made.\par
The retained coefficient checks reproduce 15 for Example 1.2, 13 for the intended homogeneous product in Example 3.6, and six for the \(2\times2\times2\) independent singular-vector specialization. Correct special cases do not resolve the repeated-vertex counterexample.\par
\section*{Conclusion}
Checked claim | Supported assessment\par
General dimension/emptiness theorem | Major: explicit admissible counterexample\par
Stated Bertini theorem | Major: counterexamples to two assertions\par
Global-section bijectivity | Minor: local false proposition\par
Connected-parameter homology claim | Minor: realizable counterexamples\par
Unconditional update convergence | Minor: divergent period-two example\par
Inspected inverse-problem arguments | No issue established within scope\par
The overall assessment is Major because the general theorem and a central proof ingredient fail under their stated assumptions. This does not imply failure of all special cases or of the inverse algorithm. No complete repaired hyperquiver theorem has been verified. Scientific decision is pending; human reviews are empty; author adoption is unknown; independent acceptance has not been performed.\par
\section*{Supporting checks}
`55f63f2ba6772c898d71d3f2f56778b910ff4`
`b34a3cbd8db5606b4b0e2461c65`.\par
The per-paper source manifest records the acquired version and source search limits.\par
The independent symbolic records are recovered unchanged from the original attempt. Its first script stopped at a structural-expression equality assertion; the preserved corrected script compares the expanded difference with zero and produced the retained successful output. This check-program failure is not a paper finding. These old calculations are not attributed to a new execution. No human review or independently verified backend equivalence is claimed. Build, dependency, parsing and report-page inspection facts are stored separately from the scientific narrative.\par
2 T. Muller, V. Nanda and A. Seigal, Multilinear Hyperquiver Representations, Foundations of Computational Mathematics, author-hosted publisher-typeset online-first version associated with 14 February 2025, https://doi.org/10.1007/s10208-025-09692-z. S. Friedland and G. Ottaviani, The number of singular vector tuples and uniqueness of best rank one approximation of tensors, consulted version arXiv:1210.8316v3, https://arxiv.org/abs/1210.8316v3.\par
\end{document}
